\chapter*{Conclusion generale et perspectives}
\addcontentsline{toc}{chapter}{Conclusion generale et perspectives}
\markboth{Conclusion generale et perspectives}{Conclusion generale et perspectives}

\section*{Bilan}
\addcontentsline{toc}{section}{Bilan}

Ce projet de fin d'annee a permis de concevoir et de realiser une plateforme
qui transforme deux sources de donnees fragmentees -- exports Qualys et
CMDB -- en un score de risque unique, contextualise et gouverne par des delais de
remediation explicites. Le travail a couvert l'ensemble de la chaine, de l'ingestion
asynchrone par lots des fichiers XML jusqu'a la restitution temps reel sous forme de
tableau de bord, en passant par la formalisation et la demonstration mathematique
complete du modele de risque (\cref{chap:modelisation}) : bornage et monotonie de
l'ARS, coherence de la classification en cinq classes, comportement du MTTR pondere
comme moyenne convexe verifiable, et modelisation du facteur de penalite operationnelle
$\mathrm{Global\_I_{Ops}}$ de l'Enterprise Risk Score. Chacune de ces proprietes a
ete verifiee numeriquement au \cref{chap:tests} sur des exemples complets, confirmant
la coherence entre la specification, sa preuve mathematique et son implementation.

Au-dela du moteur de calcul, le projet a egalement permis d'explorer des
problematiques d'ingenierie logicielle a part entiere : traitement par lots
robuste et idempotent avec Spring Batch, coherence transactionnelle d'une
architecture evenementielle asynchrone, restitution temps reel via Server-Sent
Events, et integration d'un assistant conversationnel fonde sur un modele de langage
execute localement plutot que sur une dependance externe.

\section*{Limites du travail realise}
\addcontentsline{toc}{section}{Limites du travail realise}

Plusieurs limites, assumees dans le cadre du perimetre defini
(\cref{sec:perimetre}), meritent d'etre explicitement mentionnees :

\begin{itemize}
  \item La chaine de securite HTTP autorise actuellement toute requete sans
  authentification (\texttt{permitAll}) : ce choix, adapte a un prototype
  fonctionnel demontrable, n'est pas compatible avec un usage en production sur des
  donnees aussi sensibles qu'un inventaire de vulnerabilites.
  \item La strategie d'ingestion \emph{wipe-and-reingest} (\cref{sec:pipeline-ingestion})
  simplifie considerablement le pipeline, mais empeche toute analyse d'evolution fine
  de l'inventaire CMDB lui-meme (ajout/retrait de CI) independamment du recalcul de
  risque.
  \item Le facteur $\mathrm{Global\_I_{Ops}}$ de l'ERS, entierement specifie et
  demontre au \cref{chap:modelisation}, n'est pas encore cable dans
  \texttt{DashboardMetricsService} : l'indicateur affiche est aujourd'hui une simple
  moyenne des ARS (\cref{sec:ecart-ers}).
  \item L'assistant conversationnel (RAG), decrit au \cref{chap:realisation}, reste
  au stade du developpement et de la conception : c'est un prototype fonctionnel, non
  deploye en environnement de recette (\emph{staging}) ni valide par une campagne de
  tests comparable a celle menee sur le moteur de risque (\cref{chap:tests}). Sa
  latence, identifiee des la phase de cadrage comme un point de vigilance, n'a en
  particulier pas fait l'objet d'un diagnostic instrumente permettant de determiner
  precisement la part imputable au modele, aux outils invoques ou aux acces base de
  donnees.
\end{itemize}

\section*{Perspectives}
\label{sec:perspectives}
\addcontentsline{toc}{section}{Perspectives}

Plusieurs axes d'evolution decoulent directement du bilan precedent et du perimetre
initialement ecarte :

\begin{enumerate}
  \item \textbf{Completer l'agregation ERS} en cablant le facteur
  $\mathrm{Global\_I_{Ops}}$ (\cref{eq:iops}) : ceci suppose d'introduire une cible
  SLA globale configurable et de comparer le $\mathrm{MTTR}_{global}$ mesure a cette
  cible a chaque recalcul, pour exposer aux comites de securite l'indicateur complet
  specifie.
  \item \textbf{Introduire l'authentification et le controle d'acces par role
  (RBAC)}, distinguant a minima le role Administrateur (import, configuration) du
  role Analyste (consultation), en coherence avec la generalisation d'acteurs deja
  identifiee au \cref{chap:analyse}.
  \item \textbf{Historiser la CMDB} plutot que la remplacer integralement a chaque
  import, pour permettre une analyse de tendance sur l'inventaire des actifs
  eux-memes (croissance du parc, disparition d'actifs non decommissionnes
  proprement), en complement de l'historique d'ARS deja assure par
  \texttt{RiskSnapshot}.
  \item \textbf{Renforcer la couverture de tests automatises} du moteur de risque,
  en particulier des cas limites identifies par les preuves du
  \cref{chap:modelisation} (saturation de l'ARS, absence de vulnerabilites, absence
  de correspondance CMDB, classes de risque aux bornes exactes des seuils).
  \item \textbf{Optimiser la latence de l'assistant RAG}, notamment via un cache
  (Redis) des descriptions de CVE reutilisees d'une session a l'autre -- piste deja
  identifiee lors du cadrage initial -- et via une instrumentation permettant de
  distinguer le temps passe dans le modele, dans les outils et dans les acces base
  de donnees.
  \item \textbf{Notifier proactivement} les equipes proprietaires (email, messagerie
  d'entreprise) des lors qu'un actif franchit le seuil de violation SLA
  (\cref{eq:sla-status}), plutot que de s'appuyer uniquement sur la consultation
  active du tableau de bord.
  \item \textbf{Elargir les sources ingerees} au-dela de Qualys et d'un format de
  CMDB unique, pour couvrir d'autres scanners de vulnerabilites et systemes de
  gestion de configuration frequemment rencontres chez les clients de \organisme.
  \item \textbf{Produire des exports de restitution} (rapport PDF ou tableur) a
  destination des comites de securite periodiques, en particulier pour la revue des
  actifs places en exception (\cref{eq:sla-status}), explicitement prevue par la
  gouvernance mais non encore outillee.
\end{enumerate}

Ces perspectives ne remettent pas en cause le coeur du modele mathematique
demontre dans ce rapport : elles en etendent la couverture operationnelle et
consolident sa mise en production, dans la continuite directe du travail realise
durant ce stage.
